% ============================================================
% KIPS 학술발표대회 논문 양식 (2~4페이지)
% 양식 참조: KIPS 학술발표대회 논문양식_v1_word.doc
% 규정: https://ask.kips.or.kr/notice/article/1294
% ============================================================
\documentclass[10pt, twocolumn]{article}

% --- 한국어 지원 ---
\usepackage[utf8]{inputenc}
\usepackage{kotex}

% --- 수식 ---
\usepackage{amsmath, amssymb, amsfonts}

% --- 표, 그림 ---
\usepackage{graphicx}
\usepackage{booktabs}
\usepackage{multirow}
\usepackage{caption}

% --- 알고리즘 ---
\usepackage{algorithm}
\usepackage{algpseudocode}

% --- 하이퍼링크 ---
\usepackage[hidelinks]{hyperref}

% --- 페이지 설정 (KIPS 양식 기준) ---
\usepackage{geometry}
\geometry{a4paper, top=30mm, bottom=25mm, left=20mm, right=20mm, columnsep=8mm}

% --- 제목/저자 스타일 ---
\usepackage{titling}
\setlength{\droptitle}{-2em}

% --- 헤더/푸터 ---
\usepackage{fancyhdr}
\pagestyle{fancy}
\fancyhf{}
\renewcommand{\headrulewidth}{0pt}
\fancyfoot[C]{\footnotesize 한국정보처리학회 학술발표대회 논문집}

% --- 섹션 스타일 ---
\usepackage{titlesec}
\titleformat{\section}{\normalfont\bfseries}{\thesection.}{0.5em}{}
\titleformat{\subsection}{\normalfont\bfseries}{\thesubsection}{0.5em}{}
\titlespacing*{\section}{0pt}{1.5ex}{1ex}
\titlespacing*{\subsection}{0pt}{1ex}{0.5ex}

\begin{document}

% ============================================================
% 제목 (한국어 + 영어)
% ============================================================
\twocolumn[
\begin{center}
{\Large\bfseries 강화학습 기반 5G 네트워크 슬라이싱의\\Three-Part Tariff 동적 가격 최적화}

\vspace{1em}
{\normalsize 저자명\textsuperscript{1}, 공동저자\textsuperscript{2}}\\
{\small \textsuperscript{1}소속대학교 학과}\\
{\small \textsuperscript{2}소속대학교 학과}\\
{\small \texttt{email@university.ac.kr}}

\vspace{0.8em}
{\normalsize\bfseries Dynamic Pricing Optimization for 5G Network Slicing\\via Three-Part Tariff using Reinforcement Learning}

\vspace{0.5em}
{\small Author Name\textsuperscript{1}, Co-Author\textsuperscript{2}}\\
{\small \textsuperscript{1}Dept.\ of Computer Science, University}\\
{\small \textsuperscript{2}Dept.\ of Computer Science, University}

\vspace{1em}
\begin{minipage}{0.9\textwidth}
\small\textbf{요\qquad 약}\\
\small
5G 네트워크 슬라이싱 환경에서 서비스 제공자의 수익을 극대화하기 위해 Three-Part Tariff(3PT) 구조의 동적 가격 결정 문제를 마르코프 결정 과정(MDP)으로 정식화하고, 강화학습(RL)으로 해결하는 프레임워크를 제안한다. 제안 모델은 URLLC와 eMBB 두 슬라이스에 대해 고정요금 $F$와 초과요금 $p$를 매 시간 동적으로 조정하며, 환경 모델은 로그정규 트래픽 분포, 로지스틱 이탈 모델, 가격 탄력적 Poisson 유입, 그리고 3GPP TS~22.261 기반 QoS 패널티 등 기존 문헌에 근거하여 구성된다. Soft Actor-Critic(SAC)과 Proximal Policy Optimization(PPO) 두 RL 알고리즘을 4종 베이스라인과 비교한 결과, PPO와 SAC가 각각 정적 가격 정책 대비 162.9\%, 130.2\% 높은 순보상을 달성하여 동적 가격 정책의 우수성을 확인하였다.
\end{minipage}
\vspace{1.5em}
\end{center}
]

% ============================================================
% 1. 서론
% ============================================================
\section{서론}

5G 네트워크 슬라이싱은 단일 물리적 인프라를 다수의 논리적 네트워크로 분할하여 차별화된 서비스 품질(QoS)을 보장하는 핵심 기술이다[1].
URLLC(Ultra-Reliable Low-Latency Communication) 슬라이스는 1\,ms 이내 지연과 99.999\% 신뢰성을, eMBB(enhanced Mobile Broadband) 슬라이스는 100\,Mbps 이상 처리량을 목표로 한다[2].
5G 슬라이싱 시장은 2025년 8.4억 달러에서 2030년 50.7억 달러로 연평균 43.3\% 성장이 전망된다[3].

이질적 슬라이스에 대한 \textbf{가격 설계}는 서비스 제공자의 수익과 가입자 유지에 직결되는 핵심 과제이다.
실무에서 통신 요금제는 주로 \textbf{Three-Part Tariff(3PT)}---고정 접속료~$F$, 무료 허용량~$\bar{Q}$, 초과 단위당 요금~$p$---구조를 따르며[4], 이는 Two-Part Tariff(2PT)이나 정액제보다 비용 효율적임이 이론적으로 입증되었다[4].
그러나 기존 3PT 최적화 연구[4]는 정적 수요를 가정하여, 가입자 이탈·유입의 동적 변화와 확률적 QoS 변동을 고려하지 못하는 한계가 있다.

강화학습(RL)은 모델 프리 방식으로 확률적 환경에서 연속 액션 공간의 장기 목표를 최적화할 수 있어 이 문제에 적합하다.
Kim et al.[1]은 Q-학습으로 네트워크 슬라이싱 자원 할당을 다루었으나, 종단 사용자 요금제 설계는 다루지 않았다.

본 연구의 기여는 다음과 같다:
\begin{itemize}
  \item \textbf{C1}: URLLC/eMBB 다중 슬라이스에 대한 3PT 동적 가격 결정 문제의 MDP 정식화
  \item \textbf{C2}: 로그정규 트래픽[5], 로지스틱 이탈[6,7], Poisson 유입[8], QoS 패널티[1] 등 문헌 근거 환경 모델 구축
  \item \textbf{C3}: SAC 및 PPO와 다수 베이스라인 간 실증 비교를 통한 동적 가격 정책의 우수성 검증
\end{itemize}

% ============================================================
% 2. 시스템 모델
% ============================================================
\section{시스템 모델}

\subsection{Three-Part Tariff 청구 구조}

슬라이스 $s \in \{U, E\}$에서 사용자 $i$의 시각 $t$ 청구액은[4]:
\begin{equation}
  B_{i,s,t} = F_{s,t} + \max\!\big(0,\; q_{i,s,t} - \bar{Q}_s\big) \cdot p_{s,t}
  \label{eq:bill}
\end{equation}
여기서 $F_{s,t}$는 고정요금(학습 대상), $\bar{Q}_s$는 허용량(고정), $p_{s,t}$는 초과요금(학습 대상)이다.

\subsection{데이터 사용량 모델}

개별 사용자의 시간당 데이터 사용량은 로그정규분포를 따른다[5]:
\begin{equation}
  q_{i,s,t} \sim \text{LogNormal}(\mu_s, \sigma_s^2)
  \label{eq:usage}
\end{equation}
Alasmar et al.[5]의 장기 실증 연구에 의해 인터넷 트래픽이 가우시안이 아닌 로그정규분포를 따름이 확인되었다.

\subsection{가입자 이탈 모델}

기존 가입자의 이탈 확률은 로지스틱 함수로 모델링한다[6,7]:
\begin{equation}
  P_{s,t}^{\text{dep}} = \sigma\!\Big(\gamma_{0,s} + \gamma_F \tilde{F}_{s,t} + \gamma_p \tilde{p}_{s,t} - \gamma_{\eta,s}\, \eta_{s,t-1}\Big)
  \label{eq:depart}
\end{equation}
여기서 $\sigma(x) = 1/(1+e^{-x})$, $\tilde{F}_{s,t} = F_{s,t}/F_{\text{ref},s}$는 정규화 가격, $\eta_{s,t-1}$은 이전 시간 QoS이다. $\gamma_F, \gamma_p > 0$이므로 높은 가격은 이탈을 증가시키고, $\gamma_{\eta,s} > 0$이므로 양호한 QoS는 이탈을 감소시킨다. 특히 URLLC의 QoS 민감도가 eMBB의 6배로 설정된다($\gamma_{\eta,U} \gg \gamma_{\eta,E}$). 각 가입자 $i \in \{1, \ldots, N_{s,t}\}$는 독립적으로 $\text{Leave}_{i,s,t} \sim \text{Bernoulli}(P_{s,t}^{\text{dep}})$를 따르며, 이탈자 수는 $N_{s,t}^{\text{leave}} = \sum_i \text{Leave}_{i,s,t}$이다.

\subsection{가입자 유입 모델}

신규 가입 확률은 현재 가격에만 의존한다[8, 10] (QoS 경험 없음):
\begin{equation}
  P_{s,t}^{\text{arr}} = \sigma\!\Big(\beta_{0,s} - \beta_F \tilde{F}_{s,t} - \beta_p \tilde{p}_{s,t}\Big)
  \label{eq:arrive}
\end{equation}
신규 가입자 수는 $N_{s,t}^{\text{new}} \sim \text{Poisson}(\lambda_{\max,s} \cdot P_{s,t}^{\text{arr}})$이다. 잔존자 수는 $N_{s,t}^{\text{surv}} = N_{s,t} - N_{s,t}^{\text{leave}}$이며, 시각 $t$의 활성 사용자 집합은 $\mathcal{A}_{s,t} = \{\text{잔존자}\} \cup \{\text{신규 가입자}\}$, 그 크기는 $N_{s,t}^{\text{active}} = N_{s,t}^{\text{surv}} + N_{s,t}^{\text{new}}$로 정의된다.

\subsection{QoS 및 보상 설계}

QoS 만족도는 본 모델에서 \textbf{외생적} 확률 변수 $\eta_{s,t} \sim \text{Uniform}(\eta_s^{\text{low}}, \eta_s^{\text{high}})$로 모델링한다. 즉, 에이전트는 가격을 통해 QoS를 직접 제어할 수 없으며, 관측된 QoS 변동에 대해 가격을 \emph{반응적}으로 조정하여 이탈을 완화한다(자원 할당과 가격 결정의 결합 최적화는 향후 연구).
URLLC의 목표 QoS는 3GPP TS~22.261[2]에 따라 $\eta_U^{\text{tgt}} = 0.99999$, eMBB는 $\eta_E^{\text{tgt}} = 0.90$으로 설정한다.
보상은 총 수익에서 SLA 위반 패널티를 감산한다[1]:
\begin{equation}
  r_t = \underbrace{\sum_{s} \sum_{i \in \mathcal{A}_{s,t}} B_{i,s,t}}_{R_t\text{: 총 수익}} - \underbrace{\sum_{s} w_s N_{s,t}^{\text{active}} \max(0, \eta_s^{\text{tgt}} - \eta_{s,t})}_{\text{Pen}_t\text{: QoS 패널티}}
  \label{eq:reward}
\end{equation}
여기서 $w_s$ [USD/사용자]는 단위 사용자당 단위 QoS 미달분에 대한 SLA 패널티 가중치이며, URLLC의 엄격한 요구를 반영하여 $w_U \gg w_E$로 설정한다($w_U = 500$, $w_E = 50$).

% ============================================================
% 3. MDP 정식화 및 강화학습
% ============================================================
\section{MDP 정식화 및 강화학습}

\subsection{MDP 구성}

\begin{table}[h]
\centering
\caption{MDP 정의}
\label{tab:mdp}
\small
\begin{tabular}{@{}ll@{}}
\toprule
구성요소 & 정의 \\
\midrule
상태 & $\mathcal{S}_t = (N_{U,t},\, N_{E,t},\, \eta_{U,t-1},\, \eta_{E,t-1})$ \\
액션 & $a_t = (F_{U,t},\, p_{U,t},\, F_{E,t},\, p_{E,t})$ \\
보상 & $r_t = R_t - \text{Pen}_t$ \\
전이 & $N_{s,t+1} = N_{s,t}^{\text{active}}$, $\mathcal{S}_{t+1} = (N_{U,t+1}, N_{E,t+1}, \eta_{U,t}, \eta_{E,t})$ \\
초기 & $\mathcal{S}_1 = (N_{U,1}, N_{E,1}, \eta_{U,0}, \eta_{E,0})$ \\
목표 & $\max_\pi \mathbb{E}\!\big[\sum_{t=1}^{T} \gamma^{t-1} r_t\big],\; \gamma \in [0,1)$ \\
\bottomrule
\end{tabular}
\end{table}

상태는 4차원이며, 가입자 수는 초기값으로 정규화하여 스케일 균형을 맞춘다($N_{s,t}/N_{s,1}$). 액션은 4차원 연속값(슬라이스별 고정요금 + 초과요금)이다. 에이전트는 QoS를 직접 제어할 수 없으나, QoS 변동을 관찰하여 반응적 가격 조정이 가능하다. 학습 안정성을 위해 보상에 $10^{-5}$ 스케일링을 적용한다.

\subsection{SAC와 PPO}

연속 4차원 액션 공간을 다루기 위해 두 가지 대표적 심층 강화학습 알고리즘을 적용한다. \textbf{SAC} (Soft Actor-Critic)[9]는 off-policy 방식으로, 엔트로피 정규화된 목적함수 $\max_\phi \mathbb{E}[Q_\theta(s,a) - \alpha \log \pi_\phi(a|s)]$를 최대화하며 twin Q-네트워크로 과대평가를 방지하고 자동 온도 조정($\alpha$)으로 탐색-활용 균형을 조율한다. \textbf{PPO} (Proximal Policy Optimization)는 on-policy 방식으로, 클립된 대리 목적함수를 통해 정책 갱신의 안정성을 보장한다. 두 알고리즘 모두 본 연구에서 제안하는 동적 가격 정책의 후보로 평가된다.

% ============================================================
% 4. 실험 및 결과
% ============================================================
\section{실험 및 결과}

\subsection{실험 설정}

환경 파라미터는 표~\ref{tab:params}와 같이 설정하였다. 모든 파라미터는 기존 문헌에 근거하여 캘리브레이션하였으며, 기준 가격에서 사용자별 시간당 이탈 확률 $P^{\text{dep}} = \sigma(-10.77) \approx 2.10 \times 10^{-5}$를 만족하도록 $\gamma_{0,s}$를 조정하였다(720시간 누적 기준 단순 이탈률 약 1.5\%). 정상상태 가입자 수는 유입($\lambda_{\max,s} \cdot P^{\text{arr}}$)과 이탈($N_s \cdot P^{\text{dep}}$)의 균형으로 결정된다.

\begin{table}[h]
\centering
\caption{주요 환경 파라미터}
\label{tab:params}
\small
\begin{tabular}{@{}llcc@{}}
\toprule
범주 & 파라미터 & URLLC & eMBB \\
\midrule
트래픽 & $\mu_s$ / $\sigma_s^2$ & 1.0 / 0.5 & 3.0 / 0.8 \\
       & $\mathbb{E}[q]$ (GB/hr) & 3.49 & 29.96 \\
       & $\bar{Q}_s$ (GB/hr) & 5 & 30 \\
\midrule
이탈 & $\gamma_{0,s}$ & $-9.72$ & $-12.12$ \\
     & $\gamma_{\eta,s}$ & 3.0 & 0.5 \\
\midrule
유입 & $\lambda_{\max,s}$ & 0.05 & 0.15 \\
\midrule
QoS & $\eta_s^{\text{tgt}}$ & 0.99999 & 0.90 \\
    & $w_s$ & 500 & 50 \\
\midrule
MDP & $T$ / $\gamma$ & \multicolumn{2}{c}{720 / 0.99} \\
    & $N_{s,1}$ & 1{,}000 & 5{,}000 \\
\bottomrule
\end{tabular}
\end{table}

RL 알고리즘은 Stable-Baselines3으로 구현하였으며, Actor/Critic 모두 ReLU 활성화의 3층 MLP(256-256-256)를 사용하였다. 공통 학습률은 $3 \times 10^{-4}$, 할인율 $\gamma = 0.99$이며, SAC는 배치 크기 256, 리플레이 버퍼 $10^6$, 소프트 업데이트 $\tau = 0.005$, 자동 엔트로피 조정($\alpha=$auto)을, PPO는 배치 크기 64, $n_{\text{epochs}}=10$, 클립 범위 $0.2$, GAE $\lambda = 0.95$를 적용하였다. 가입자 수 관측값은 초기값으로 정규화하고, 보상은 $10^{-5}$ 스케일링하여 학습 안정성을 확보하였다. SAC는 500 에피소드, PPO는 1{,}000 에피소드 학습 후, 20회 평가 에피소드의 평균으로 성능을 측정하였다(시드 42, 단일).

\subsection{베이스라인}

학습된 RL 정책의 우수성을 검증하기 위해 다음 베이스라인과 비교한다:
(1)~\textbf{Static-Heuristic}: 기준가격($F_{\text{ref}}, p_{\text{ref}}$) 고정,
(2)~\textbf{Random}: 균등 랜덤 액션,
(3)~\textbf{2PT-SAC}: 허용량 $\bar{Q}=0$으로 SAC 학습,
(4)~\textbf{Flat-SAC}: 초과요금 $p=0$으로 고정하여 $F$만 학습.

\subsection{실험 결과}

표~\ref{tab:results}는 각 정책의 성능을 비교한 결과이다.

\begin{table}[h]
\centering
\caption{정책별 성능 비교 (20회 평가, 평균$\pm$표준편차)}
\label{tab:results}
\scriptsize
\begin{tabular}{@{}lrrr@{}}
\toprule
정책 & 순보상\footnotemark & 수익 (M\$) & 패널티 (M\$) \\
\midrule
PPO (제안) & $\mathbf{8{,}194 \pm 38}$ & $\mathbf{840.8 \pm 3.7}$ & $21.41 \pm 0.41$ \\
SAC (제안) & $7{,}176 \pm 64$ & $738.5 \pm 6.4$ & $\mathbf{20.84 \pm 0.34}$ \\
Max-Price & $5{,}870 \pm 48$ & $606.3 \pm 4.8$ & $19.30 \pm 0.52$ \\
Random & $5{,}342 \pm 48$ & $556.4 \pm 4.9$ & $22.24 \pm 0.47$ \\
Static-Heur. & $3{,}117 \pm 5$ & $334.3 \pm 0.5$ & $22.60 \pm 0.40$ \\
Zero-Price & $-228 \pm 4$ & $0$ & $22.80 \pm 0.40$ \\
\bottomrule
\end{tabular}
\end{table}
\footnotetext{순보상은 $10^{-5}$ 스케일링 적용 값이며, 수익/패널티는 원시 USD 값으로 보고한다. PPO는 전체 정책 중 최고 순보상/수익을, SAC는 RL 정책 중 최저 패널티를 달성하여 각각 굵게 표시하였다.}

Random 정책이 Static-Heuristic보다 높은 순보상을 달성하는 것은, 무작위 탐색이 기준가격보다 높은 수익 구간을 빈번히 탐색하기 때문이다. 그러나 Random 정책은 가입자 수가 불안정하게 변동하며($N_E$: 5{,}000 $\rightarrow$ 4{,}358), 장기적으로 지속 가능하지 않다.
PPO(순보상 8{,}194)와 SAC(7{,}176) 모두 전체 베이스라인을 상회하였다.
PPO는 Static-Heuristic 대비 162.9\%, Max-Price 대비 39.6\% 높은 순보상을 달성하였으며, SAC 대비 14.2\% 우수하였다.
PPO가 SAC를 상회한 것은 on-policy 방식이 이 환경의 비정상성(가입자 수 점진적 변화)에 더 적응적으로 반응했기 때문으로 분석된다.
한편, SAC는 패널티 20.8M으로 PPO(21.4M) 및 모든 정적 베이스라인 대비 낮은 수준을 유지하여 RL 정책 중 QoS 준수 측면에서 가장 우수하였다. Max-Price(19.3M)가 더 낮은 패널티를 보이는 것은 가입자 급감으로 인한 사용자 기반 축소의 부작용이며, 이는 지속 가능한 정책이 아니다.
Max-Price 정책은 eMBB 가입자가 5{,}000명에서 909명으로 급감하여 장기 지속 불가능한 반면, PPO는 4{,}102명, SAC는 2{,}188명을 유지하며 가입자 기반과 수익의 균형을 학습하였다.

공유 파라미터: $\gamma_F = 1.0$, $\gamma_p = 0.8$, $\beta_{0,U} = 2.0$, $\beta_{0,E} = 2.5$, $\beta_F = 0.8$, $\beta_p = 0.6$, $F_{\text{ref}} = (50, 30)$, $p_{\text{ref}} = (10, 5)$.


% ============================================================
% 5. 결론
% ============================================================
\section{결론}

본 연구에서는 5G 네트워크 슬라이싱 환경에서 Three-Part Tariff 가격을 강화학습으로 동적 최적화하는 MDP 프레임워크를 제안하였다.
제안 모델은 로그정규 트래픽, 로지스틱 이탈, Poisson 유입, 3GPP QoS 기준 등 문헌에 근거한 환경 위에서 SAC와 PPO 에이전트가 URLLC/eMBB 슬라이스별 고정요금과 초과요금을 동적으로 결정한다.
시뮬레이션 결과, PPO와 SAC 에이전트는 각각 Static-Heuristic 대비 162.9\%, 130.2\% 높은 순보상을 달성하였으며, 모든 고정 가격 베이스라인을 상회하여 강화학습 기반 동적 가격 정책의 우수성을 확인하였다.
본 연구의 한계로는 단일 시드(42) 결과만을 보고한 점을 들 수 있으며, 다중 시드 검증 및 통계적 유의성 분석은 향후 과제로 남긴다.
향후 연구에서는 외생적 QoS를 내생화하여 가격-자원 할당 결합 최적화, 다중 사업자 경쟁 모델링, 그리고 mMTC 슬라이스 확장을 다룰 계획이다.

% ============================================================
% 사사
% ============================================================
% \section*{사사}
% 본 연구는 ... 의 지원을 받아 수행되었습니다.

% ============================================================
% 참고문헌
% ============================================================
\begin{thebibliography}{9}

\bibitem{kim2019}
Y.~Kim, S.~Kim, H.~Lim, ``Reinforcement Learning Based Resource Management for Network Slicing,'' \textit{Applied Sciences}, Vol.~9, No.~11, p.~2361, 2019.

\bibitem{3gpp22261}
3GPP, ``Service Requirements for the 5G System,'' 3GPP TS~22.261, 2023.

\bibitem{5gmarket}
Research and Markets, ``5G Network Slicing Market Size, Share \& Trends Analysis Report,'' 2025.

\bibitem{wu2023}
J.~F. Wu, B.~Behzad, ``Optimal Three-Part Tariff Pricing with Spence-Mirrlees Reservation Prices,'' \textit{Math. Methods Oper. Res.}, Vol.~97, pp.~233--258, 2023.

\bibitem{alasmar2021}
M.~Alasmar, R.~Clegg, N.~Zakhleniuk, G.~Pavlou, ``Internet Traffic Volumes Are Not Gaussian---They Are Log-Normal,'' \textit{IEEE/ACM Trans. Netw.}, Vol.~29, No.~3, pp.~1266--1279, 2021.

\bibitem{ahn2006}
J.-H. Ahn, S.-P. Han, Y.-S. Lee, ``Customer Churn Analysis: Churn Determinants and Mediation Effects of Partial Defection in the Korean Mobile Telecommunications Service Industry,'' \textit{Telecom. Policy}, Vol.~30, pp.~552--568, 2006.

\bibitem{ribeiro2023}
H.~Ribeiro, B.~Barbosa, A.~C.~Moreira, R.~G.~Rodrigues, ``Determinants of Churn in Telecommunication Services: A Systematic Literature Review,'' \textit{Mgmt Rev. Quarterly}, Vol.~74, pp.~1327--1364, 2023.

\bibitem{chen2022}
J.~Chen, L.~Jiang, S.~A.~S.~Shah, ``An Empirical Model of the Effects of `Bill Shock' Regulation in Mobile Telecommunication Markets,'' \textit{Int. J. Ind. Organ.}, Vol.~83, p.~102848, 2022.

\bibitem{haarnoja2018}
T.~Haarnoja, A.~Zhou, P.~Abbeel, S.~Levine, ``Soft Actor-Critic: Off-Policy Maximum Entropy Deep Reinforcement Learning with a Stochastic Actor,'' in \textit{Proc. ICML}, 2018.

\bibitem{jiang2014}
L.~Jiang, ``An Empirical Model of the Effects of `Bill Shock' Regulation in Mobile Telecommunication Markets,'' Working Paper, Univ.\ of California, Irvine, 2014.

\end{thebibliography}

\end{document}
